%% 2. RELATED WORK
%% ============================================================================

%% 2.1 LLM-Based Code Generation

The field of LLM-based code generation has seen rapid advancement. GitHub Copilot demonstrated the commercial viability of AI-assisted coding, while subsequent research has focused on improving code quality, reducing hallucinations, and integrating with development workflows.

Recent work has explored the use of larger context windows to enable whole-repository reasoning. However, as noted by Liu et al. (2024), long context alone does not solve the structural reasoning problem - the model still lacks explicit representations of code structure.

%% 2.2 Structural Analysis in Software Engineering

Traditional software engineering tools have long used structural analysis for quality assurance:

\begin{itemize}
    \item STATIC ANALYSIS: Tools like SonarQube, PMD, and Checkstyle analyze ASTs to detect code smells and anti-patterns.
    \item TYPE SYSTEMS: Modern type systems (TypeScript, Python type hints, Rust) catch errors at compile time by understanding code structure.
    \item LINTERS: ESLint, Pylint, and similar tools enforce coding standards using AST traversal.
\end{itemize}

%% 2.3 Code Representation Learning

Learning representations of code structure has been an active research area:

\begin{itemize}
    \item CODEBERT (Feng et al., 2020): A pre-trained model for programming and natural languages using a multi-modal transformer architecture.
    \item GRAPHCODEBERT (Guo et al., 2020): Incorporates data flow into pre-training for better code understanding.
    \item CODE2VEC (Alon & Shalev, 2019): Learns distributed representations of code using AST paths.
    \item PLBART (Ahmad et al., 2021): A pre-trained model for program and natural language that can handle code generation and understanding.
\end{itemize}

%% 2.4 Hybrid LLM-Specialist Architectures

Our work builds on the emerging paradigm of hybrid architectures:

\begin{itemize}
    \item TOOL-USING AGENTS: Systems that call external tools for specific tasks (Schick et al., 2023; Qin et al., 2023).
    \item MIXTURE OF EXPERTS: Models that route different types of problems to specialized sub-models (Shazeer et al., 2017; Fedus et al., 2021).
    \item RETRIEVAL-AUGMENTED GENERATION: Systems that retrieve relevant information to augment LLM outputs (Lewis et al., 2020; Guu et al., 2020).
\end{itemize}

%% 2.5 State Space Models for Code

State Space Models (SSMs) have gained attention for their efficiency and theoretical properties:

\begin{itemize}
    \item MAMBA (Gu & Goel, 2023): A selective SSM that achieves linear scaling with respect to sequence length while maintaining competitive performance with Transformers.
    \item MAMBA-2 (Gu et al., 2024): Improves upon the original Mamba architecture with better performance and efficiency.
    \item APPLICATIONS TO CODE: Recent work has explored using SSMs for code tasks due to their efficiency and ability to handle long sequences.
\end{itemize}

Our work is the first to combine SSM-based embeddings with structural analysis for real-time code quality gating in LLM-based development workflows.

%% 2.6 Evaluation of Code Generation

Evaluating code generation systems presents unique challenges:

\begin{itemize}
    \item HUMANEVAL (Chen et al., 2021): A benchmark for evaluating functional correctness of code generation.
    \item MBPP (Austin et al., 2021): A benchmark with 1,000 Python problems for assessing code generation.
    \item SWE-BENCH (Jin et al., 2024): A benchmark for evaluating LLMs on real-world software engineering tasks.
    \item AGENTBENCH (Liu et al., 2023): A comprehensive benchmark for LLM-based agents.
\end{itemize}

Our evaluation methodology extends these approaches by incorporating structural quality metrics and multi-judge cross-family assessment.

%% PLACEHOLDER: Add comparison to specific papers on LLM code review
%% PLACEHOLDER: Add discussion of concurrent work
%% PLACEHOLDER: Add positioning relative to commercial products
